% 27-ceresmip-tier-2 — CERESMIP, "Tier 2": the ten Tier 2 experiments as
% renumbered on the GISS site (2026-02-02), all 1990--2024, and the CESM3/CAM7
% runs planned against them. Science-level only. Facts:
% ../CERES-MIP/fable/02-ceresmip-protocol.md (Table 2) and 09-run-plan.md.
% Layout budget as slide 02: ten body lines in the right column (two blocks
% of three items, one of them single-line), lines under ~34 characters in
% this narrower column. The table column has room to spare.
\begin{frame}[t]
\frametitle{Tier 2}
\footnotesize
\begin{columns}[T]
\begin{column}{0.51\textwidth}
\centering
\textbf{Tier 2 Experiments, 1990--2024}\\[2pt]
{\scriptsize
\begin{tabular}{@{}ll@{}}
\toprule
Experiment & Design \\
\midrule
amip ensemble & all forcings, 3+ members \\
amip-nudge & winds nudged to reanalysis \\
amip-nat & solar, volcanic, orbital only \\
amip-ghg & greenhouse gases only \\
amip-aer & anthropogenic aerosol only \\
amip-slcf & short-lived forcers only \\
amip-sst & SST and sea-ice change only \\
amip-noghg & all but greenhouse gases \\
amip-noaer & all but aerosol \\
piclim-histall & all forcings, 1850 SST \\
\bottomrule
\end{tabular}\\[4pt]
Single-forcing runs: other inputs at 1850.\par}
\end{column}
\begin{column}{0.47\textwidth}
\begin{block}{Planned CESM3/CAM7 Runs}
\begin{itemize}
  \item Ensembles of 3--5 members,\\free and AOD-constrained
  \item Aerosol-only and all-but-aerosol
  \item Nudged winds: ERA5,\\as in the 2008 test run
\end{itemize}
\end{block}
\begin{block}{What the Pairs Separate}
\begin{itemize}
  \item Free vs constrained:\\aerosol state or aerosol response
  \item Ensemble: internal variability
  \item All-VIIRS member, 2018--2024:\\the April 2022 imager seam
\end{itemize}
\end{block}
\end{column}
\end{columns}
\end{frame}
